\section{分析与可解释性 (Analysis and Interpretability)}
\label{sec:analysis}

\subsection{章节定位 + 零基础该问的问题}

前面几章我们一路堆模型：先有 seq2seq，再加 attention，再加 multi-head，
每一步都伴随着一个“直觉理由”——比如“注意力应该能学到对齐”，
“多头应该能学到不同种类的关系”。
但这些都只是\emph{设计动机}，并不等于模型真的照做了。
这一章要做的事是：\textbf{打开训练好的模型，看看里面到底发生了什么}。

\begin{problemblock}{零基础读者会问的三个问题}
\begin{enumerate}
  \item 多头注意力的每个 head 真的在干不同的事吗？还是大家学的都差不多？
  \item 既然这样，能不能干脆把一些“没用”的 head 裁掉？
  \item Encoder 内部的向量表示，到底\emph{编码了什么语言学信息}（词性？形态？句法？）？
\end{enumerate}
\end{problemblock}

下文以两条研究线分别回应：第一条来自
Voita et al.\ (ACL 2019, \emph{Analyzing Multi-Head Self-Attention}),
第二条来自 Belinkov et al.\ (2017) 等围绕 NMT encoder 的 probing 工作。

\subsection{多头自注意力的头到底在干嘛？}
\label{sec:head-roles}

\subsubsection{重要的 head 是可解释的}

Voita et al.\ 训练好一个 Transformer NMT 模型后，逐个 head 分析其行为，
得到一个非常干净的结论：\textbf{真正“干活”的只是少数 head}，
而这些 head 的角色\emph{可以被人读懂}。

\begin{intuitionblock}
不是每个 head 都是平等的。一个训练好的多头自注意力，
更像是“几个专家 + 一堆陪跑”。
\end{intuitionblock}

被识别出的“专家”大体可以归成三类：

\begin{description}
  \item[(a) 位置型 head (positional heads)]
    几乎只看相邻位置。
    Encoder 中典型地有 2--3 个“看前一个 token”的 head，
    以及 \(\sim\)2 个“看后一个 token”的 head。
    这相当于模型自己学出了一个轻量级的局部卷积。
  \item[(b) 句法型 head (syntactic heads)]
    会跟踪“主语\(\to\)动词”“动词\(\to\)宾语”这类语法关系。
    这种行为在多个不同语言对的模型里都能观察到，
    说明它不是某种偶然，而是 multi-head 自然倾向于做的事。
  \item[(c) 稀有词型 head (rare-tokens head)]
    在第一层往往会出现一个特殊的 head：它把注意力集中到句子里最稀有的 token 上。
    一个合理的解释是，模型在“锚定”那些信息量最大的稀有线索；
    同时这个现象也带有一点轻微的过拟合气味
    （稀有 token 在训练集里本来就更具区分度）。
\end{description}

这个结果反过来验证了 multi-head 的归纳偏置：
\textbf{不同 head 确实会学到不同的东西}，并不是简单地把同一种关系算了 \(h\) 遍。

\subsubsection{大多数 head 可以被剪掉}

同一篇论文进一步做了一件激进的事：
给每个 head 套上一个可学习的“门”（gate），
通过稀疏正则把不重要的 head 推向 0，
然后看模型质量随被裁 head 数量的下降。

观察到两个互相印证的现象：

\begin{enumerate}
  \item \textbf{很大一部分 head 可以被裁掉}，翻译质量几乎不掉。
        典型设置下，只剩下 1/4 \(\sim\) 1/3 的 head，BLEU 仍然接近原模型。
  \item \textbf{留下来的 head 恰好就是上面归类的那些“专家”}
        （位置型、句法型、稀有词型）。
\end{enumerate}

但有一个看起来反直觉、却非常重要的细节：

\begin{keypointblock}
\textbf{这\emph{不}意味着可以一开始就训练一个 head 数很少的模型。}
直接用少头训练，质量明显更差。
多头在训练时的真正作用更像是\emph{一个更大的假设空间}：
让模型有机会“尝试”各种行为模式，最后挑出有用的那几个。
剪枝是事后压缩，而不是先验设计。
\end{keypointblock}

\subsection{Probing：表示里到底编码了什么？}
\label{sec:probing}

第二个问题指向 encoder 输出的那一堆向量：它们里面除了“能用来翻译”，
是否还藏着具体的语言学信息？衡量这件事的标准工具叫 \textbf{probing}（探针）。

\subsubsection{Probing 的方法论}

\begin{mathbox}{Probing 的标准流程}
\begin{enumerate}
  \item \textbf{冻结}已经训练好的 NMT encoder（参数不再更新）。
  \item 把一批带标注的数据（例如每个词都有词性标签）喂进去，
        在指定层取出每个词的向量表示 \(h_i\)。
  \item 在这些\emph{冻结}的 \(h_i\) 之上，训练一个\emph{小}分类器
        （通常是线性层或浅层 MLP）去预测目标标签 \(y_i\)。
  \item 用这个小分类器的准确率作为代理指标，
        近似回答“\(h_i\) 里有多少关于 \(y_i\) 的信息”。
\end{enumerate}
\end{mathbox}

\begin{remark}
一个常被忽视的注意点：probing 的准确率本身并不直接等于
“表示里编码了多少信息”，因为探针自己就有一定学习能力，
甚至能从随机特征里学出一些东西。
严格的 probing 研究通常会做对照实验（比如随机初始化的 encoder、
打乱标签的对照、或控制探针容量）。
本节只用其结论性观察，对方法论细节不再展开。
\end{remark}

\subsubsection{NMT encoder 学到了什么形态学？}

Belinkov et al.\ (2017) 用 probing 方法系统研究了 NMT encoder。
两个值得记住的结论：

\paragraph{(a) 词性 (POS) 与层数的关系}
拿 encoder 各层的表示去 probe POS 标注：
\begin{itemize}
  \item 训练后的 encoder 各层都比原始词嵌入显著更好；
  \item 但有趣的是，\textbf{第 1 层的 POS 探针准确率往往\emph{高于}第 2 层}。
\end{itemize}
对应的解释假设是：
第 1 层更接近“词级结构”（词性、形态等表层属性），
第 2 层开始向更抽象、更语义化的方向移动，
反而把表层语法信息“压”了下去。
这与“层数越深 = 越好”的朴素直觉是矛盾的。

\paragraph{(b) 目标语言的反向影响}
固定 source 语言（比如阿拉伯语），变化 target 语言，再 probe source 端形态信息：

\begin{intuitionblock}
直觉上，target 越“强”，encoder 越要努力工作，所以 encoder 的表示应该越丰富。
\end{intuitionblock}

实际观察到的是这样一个稍微反直觉的结果：

\begin{quote}
\textbf{Target 语言形态\emph{越贫乏}（如英语），encoder 对 source 形态的编码反而\emph{越好}；
target 语言形态\emph{越丰富}（如希伯来语、德语），encoder 反而编码得更弱。}
\end{quote}

一个解释方式是：
当 target 自己有丰富的屈折变化时，模型可以“偷懒”，
让大部分形态信息直接从 target 端的解码器/词表那里产生；
而 target 一旦贫乏（英语几乎没有屈折），
encoder 就不得不在自己这一侧把 source 的形态彻底搞清楚，
否则解码器没法做出正确翻译。

\subsection{这一章真正该记住的}

\begin{keypointblock}
\begin{enumerate}
  \item \textbf{头不平等}：训练好的多头注意力里，只有一小部分 head 真正承担主要工作。
  \item \textbf{角色可解释}：重要的 head 倾向于扮演位置型、句法型、稀有词型这几种角色，
        说明 multi-head 的归纳偏置确实在起作用。
  \item \textbf{剪枝可行，但前提是先有多头}：可以事后裁掉大部分 head 而几乎不掉点；
        但直接用少头从头训练效果会明显更差——多头主要是\emph{训练阶段的假设空间}。
  \item \textbf{Probing 是“表示里有什么”的标准工具}：
        冻结 encoder，在表示之上训练浅层分类器，用准确率近似信息量；
        但要小心探针自身能力带来的偏差。
  \item \textbf{Target 越弱，encoder 越强}：
        当 target 语言形态贫乏时，NMT encoder 反而会把 source 的形态信息编码得更彻底。
\end{enumerate}
\end{keypointblock}
